% ============================================================================
% Sección 8: Reproducibilidad
% ============================================================================
\section{Reproducibilidad}

Esta sección documenta los pasos necesarios para reproducir completamente los experimentos presentados en este trabajo. Todo el código fuente es público y los resultados son deterministas.

\subsection{Repositorios del proyecto}

El ecosistema se compone de cuatro repositorios públicos:

\begin{table}[H]
\centering
\caption{Repositorios del ecosistema}
\label{tab:repos}
\renewcommand{\arraystretch}{1.3}
\begin{tabular}{@{}lll@{}}
\toprule
\textbf{Componente} & \textbf{Repositorio} & \textbf{Rol} \\
\midrule
Core library & \texttt{ahincho/search-library} & Algoritmos de búsqueda \\
Benchmark grafos & \texttt{ahincho/graph-search} & Experimentos en grafos \\
Benchmark grids & \texttt{ahincho/path-search} & Experimentos en grids \\
Artículo & \texttt{ahincho/search-paper} & Este documento \\
\bottomrule
\end{tabular}
\end{table}

Enlaces directos:
\begin{itemize}
    \item PyPI: \url{https://pypi.org/project/search-library/}
    \item Core: \url{https://github.com/ahincho/search-library}
    \item Grafos: \url{https://github.com/ahincho/graph-search}
    \item Grids: \url{https://github.com/ahincho/path-search}
    \item Paper: \url{https://github.com/ahincho/search-paper}
\end{itemize}

\subsection{Requisitos del entorno}

\begin{itemize}
    \item \textbf{Python:} $\geq$ 3.11 para \texttt{search-library}; $\geq$ 3.14 para los benchmarks.
    \item \textbf{Gestor de paquetes:} \texttt{uv} (recomendado) o \texttt{pip} tradicional.
    \item \textbf{Sistema operativo:} Multiplataforma (Windows, Linux, macOS).
    \item \textbf{Sin dependencias externas:} La librería core no requiere paquetes de terceros.
\end{itemize}

\subsection{Instalación de la librería core}

\subsubsection{Opción A: Desde PyPI (recomendado)}

\begin{lstlisting}[language=bash, caption={Instalación desde PyPI}, label={lst:install_pypi}]
# Con pip
pip install search-library

# Con uv
uv add search-library
\end{lstlisting}

\subsubsection{Opción B: Desde código fuente}

\begin{lstlisting}[language=bash, caption={Instalación desde fuente}, label={lst:install_source}]
git clone https://github.com/ahincho/search-library
cd search-library

# Con pip
pip install -e .

# Con uv
uv sync
\end{lstlisting}

\subsection{Reproducción del benchmark en grafos}

\begin{lstlisting}[language=bash, caption={Ejecución del benchmark de grafos}, label={lst:run_graph}]
git clone https://github.com/ahincho/graph-search
cd graph-search

# Opcion 1: Con uv (recomendado)
uv sync
uv run python src/main.py

# Opcion 2: Con pip tradicional
pip install search-library>=1.0.0
python src/main.py
\end{lstlisting}

\textbf{Salida esperada:} Tabla comparativa de 5 algoritmos (A*, Dijkstra, BFS, DFS, Bidirectional) sobre la red vial Arequipa $\rightarrow$ Cusco con métricas de costo, nodos explorados y optimalidad.

\subsection{Reproducción del benchmark en grids}

\begin{lstlisting}[language=bash, caption={Ejecución del benchmark de grids}, label={lst:run_grid}]
git clone https://github.com/ahincho/path-search
cd path-search

# Opcion 1: Con uv (recomendado)
uv sync
uv run grid-benchmark

# Opcion 2: Con pip tradicional
pip install search-library>=1.0.0
python src/main.py
\end{lstlisting}

\textbf{Salida esperada:} Reporte completo de 8 escenarios con tablas de comparación, visualizaciones ASCII de grids, y análisis académico automático.

\subsection{Ejecución de pruebas unitarias}

\begin{lstlisting}[language=bash, caption={Ejecución de tests en la librería core}, label={lst:tests}]
cd search-library

# Con uv
uv run pytest

# Con pip
pip install pytest pytest-cov
pytest --cov=search_library --cov-fail-under=90
\end{lstlisting}

La cobertura mínima requerida es del 90\%. El reporte incluye cobertura por línea y por rama (\texttt{branch=true}).

\subsection{Compilación del artículo}

\begin{lstlisting}[language=bash, caption={Compilación local del PDF}, label={lst:compile}]
git clone https://github.com/ahincho/search-paper
cd search-paper

# Windows (PowerShell)
.\compile.ps1

# Windows (CMD)
compile.bat

# Linux/macOS
chmod +x scripts/build.sh
./scripts/build.sh
\end{lstlisting}

El pipeline CI/CD en GitHub Actions compila automáticamente en cada push a \texttt{main} y genera releases versionados con el PDF adjunto al crear tags \texttt{v*.*.*}.

\subsection{Garantías de reproducibilidad}

\begin{itemize}
    \item \textbf{Determinismo:} Todos los grids con obstáculos aleatorios se generan con semillas explícitas (\texttt{seed}). El mismo seed produce siempre el mismo grid.
    \item \textbf{Sin estado externo:} Los benchmarks no dependen de archivos externos, bases de datos, ni servicios de red.
    \item \textbf{Versión fija:} La dependencia \texttt{search-library>=1.0.0} está publicada en PyPI con versionado semántico. Los resultados de este artículo corresponden a la versión 1.0.0.
    \item \textbf{CI/CD verificado:} El workflow de GitHub Actions compila el artículo y ejecuta los tests automáticamente, garantizando que el código en el repositorio siempre produce los resultados documentados.
\end{itemize}
